% =====================================================================
% Sección de García Olvera Emiliano
% Medición RNF2 (háptica), consumo eléctrico (Clase 16), Clase 20 y 21 (actuadores)
% =====================================================================

\subsection{Medición RNF2: Respuesta Háptica}

\textbf{Requisito RNF2:} Activación del motor $\le 50\text{ ms}$ tras señal.

\textbf{Medición pendiente:}
\begin{itemize}
	\item Tiempo medido desde comando hasta activación PWM: \pend{XX.X ms}
	\item Implementación: PWM a 20 kHz mediante alarma no bloqueante
\end{itemize}

\textbf{Veredicto:} \pend{Cumplido}

\subsection{Clase 16: Consumo Eléctrico}

\textbf{Mediciones de consumo:}
\begin{itemize}
	\item Raspberry Pi Pico: $\approx 25\text{ mA}$ a $5\text{ V}$ = $0.13\text{ W}$
	\item Raspberry Pi 4 en reposo: \pend{X.XX W}
	\item Raspberry Pi 4 bajo carga (jugando): \pend{X.XX W}
	\item Gobernador: \texttt{ondemand} (600--1800 MHz)
	\item Estado térmico: \texttt{vcgencmd get\_throttled} = \texttt{0x0}
\end{itemize}

\subsection{Clase 20 y 21: Actuadores}

\textbf{Clase 20 (Tipos):}
\begin{itemize}
	\item Comandos HID OUT de 8 bytes (tipo + parámetros)
	\item Uso de \texttt{volatile} en variables compartidas con IRQ
\end{itemize}

\textbf{Clase 21 (Drivers):}
\begin{itemize}
	\item Motor y buzzer: \texttt{hardware/pwm} (slices PWM del RP2040)
	\item LEDs: \texttt{hardware/gpio}
	\item Evidencia: configuración en \texttt{actuadores.c}
\end{itemize}
